\documentclass[11pt]{article}

\usepackage[utf8]{inputenc}
\usepackage[T1]{fontenc}
\usepackage[margin=1in]{geometry}
\usepackage[strings]{underscore} % allow _ in text / \texttt
\usepackage{xcolor}
\usepackage{listings}
\usepackage{booktabs}
\usepackage{longtable}
\usepackage{enumitem}
\usepackage{parskip}
\usepackage[hidelinks]{hyperref}

\definecolor{codebg}{rgb}{0.96,0.96,0.96}
\lstdefinestyle{mono}{
  basicstyle=\ttfamily\footnotesize,
  backgroundcolor=\color{codebg},
  breaklines=true,
  columns=fullflexible,
  keepspaces=true,
  frame=single,
  framesep=4pt,
  xleftmargin=4pt,
  showstringspaces=false
}
\lstset{style=mono}

% inline code helper. Uses \texttt (works in tables/headers, unlike \lstinline);
% the underscore package makes _ render literally. Snippets containing ^ \ % { }
% escape those characters explicitly at the call site.
\newcommand{\ic}[1]{\texttt{#1}}

\title{\textbf{Claude Code Inside Black Ops}\\[4pt]
\large Bridging the \emph{Call of Duty: Black Ops} (2010) easter-egg terminal to a
real Windows shell, and building a full terminal emulator inside the game process}
\author{Reverse-engineering \& implementation report}
\date{June 2026}

\begin{document}
\maketitle

\begin{abstract}
\noindent
The 2010 game \emph{Call of Duty: Black Ops} (engine codename \texttt{t5}, an
idTech\,3/IW derivative) contains a hidden in-game ``CIA terminal'' easter egg: a
hardcoded simulation with a fixed command table, a fake filesystem, and scripted
string responses. It has \emph{no} I/O to the host. This project bolts real
behaviour onto it entirely from \emph{outside} the binary --- with no source code,
no symbols, and no debug information --- by injecting a DLL that hooks the game's
own terminal code. Two capabilities were built on top of the same reverse
engineering: (1) a \emph{line-shell bridge} that runs typed commands in a real
persistent \texttt{cmd.exe} and renders the output through the game's own text
renderer; and (2) a \emph{full terminal emulator} --- a Windows pseudoconsole
(ConPTY) driving a VT/ANSI parser and an 80$\times$27 character grid that is blitted
straight into the engine's display memory each frame --- inside which interactive,
full-screen programs run. The headline result: \textbf{Claude Code itself runs and
renders correctly inside the Black Ops terminal}, with colour, live spinners, and
keyboard input. This report documents how every piece works in detail, the notable
discoveries made while reversing the binary, and every error encountered together
with its fix.
\end{abstract}

\tableofcontents

\section{Goal and constraints}

The objective was \emph{authenticity}: the work had to happen inside the real game
process and render in the game's own terminal, not in a look-alike standalone
window. A themed external terminal was explicitly considered and rejected.

Hard constraints that shaped every decision:

\begin{itemize}[nosep]
  \item \textbf{Offline / single-player only.} Black Ops multiplayer ran VAC
  (Valve Anti-Cheat); injecting into a VAC-protected session is a ban. The terminal
  is a campaign/menu feature, so all development and use is done with networking off.
  \item \textbf{Personal, single-machine novelty.} The bridge runs whatever is typed
  in-game as a host shell command --- arbitrary code execution \emph{by design}. It
  is never exposed over a network; any future IPC channel must bind to localhost only.
  \item \textbf{No source, no symbols.} Everything was recovered from the shipped
  \texttt{BlackOps.exe} via static analysis (Ghidra) plus live observation.
\end{itemize}

\section{The target binary}

\begin{longtable}{@{}ll@{}}
\toprule
\textbf{Property} & \textbf{Value} \\
\midrule
Path & \texttt{C:\textbackslash Games\textbackslash Call of Duty - Black Ops\textbackslash BlackOps.exe} \\
Size & 8{,}099{,}608 bytes \\
Machine & \texttt{0x014C} (i386 --- \textbf{32-bit x86}) \\
Optional header magic & \texttt{0x10B} (PE32) \\
Image base & \texttt{0x00400000} \\
ASLR & \textbf{off} (\texttt{DYNAMIC_BASE} clear) --- always loads at \texttt{0x400000} \\
\bottomrule
\end{longtable}

\noindent
Two facts here are load-bearing for the whole project. First, the binary is
\textbf{32-bit}, which dictates an x86 DLL, an x86 debugger, and --- as discovered
much later --- causes WOW64 file-system redirection that hides 64-bit-only tools
(Section~\ref{sec:ssh}). Second, \textbf{ASLR is off}: the module always loads at
\texttt{0x400000}, so a runtime virtual address equals \texttt{0x400000 + RVA}.
That makes addresses stable across runs and lets globals be reached as
\ic{module_base + RVA}. Even so, code is hooked by \emph{byte-pattern signature}
where practical, so a game update that shifts code does not silently break a
hardcoded address; absolute address operands inside an instruction are wildcarded
in the pattern.

\section{Reverse-engineering methodology}

The way into a closed binary like this is the terminal's own \textbf{hardcoded
response strings}. The recipe:

\begin{enumerate}[nosep]
  \item In Ghidra, list \emph{Defined Strings} and find a known terminal response
  (HELP output, a Zork line, a login prompt). 116 candidate strings were exported.
  \item Follow that string's \emph{cross-references}. The referencing code is the
  command-dispatch logic, or one hop from it.
  \item Set a breakpoint there in a live debugger, type the matching command in-game,
  and read the registers/stack at the break.
  \item Walk the call stack \emph{up} to where the raw input line is assembled
  (the submit handler), and \emph{down} to where a response string is handed to the
  renderer (the line-append function).
\end{enumerate}

All static recon was scripted with Ghidra \emph{headless} (\texttt{analyzeHeadless},
\texttt{-noanalysis} against the already-analysed project), so it could be re-run
cheaply. Those scripts live in \texttt{ghidra_scripts/} (e.g.
\texttt{DumpTerminalStrings}, \texttt{FindTermRender}, \texttt{ReconTermEmu},
\texttt{ReconInput}, \texttt{DumpDraw}, \texttt{DumpEmit}); their decompiled output
is captured in \texttt{notes/} and the curated address table in
\texttt{notes/offsets.md}.

\section{The engine functions that were recovered}
\label{sec:functions}

\begin{longtable}{@{}p{3.1cm}p{1.7cm}p{8.3cm}@{}}
\toprule
\textbf{Symbol} & \textbf{RVA} & \textbf{Role} \\
\midrule
\endhead
\texttt{line_append}\newline(FUN_004ef440) & \texttt{0xEF440} &
The authentic scrollback renderer. \ic{__cdecl(const char* text, int color=1, int flag=1, int zero=0)}. Parses \ic{\textasciicircum N} colour codes, wraps at 80 cols, scrolls at 21 rows, draws glyphs via \texttt{FUN_004165b0}. \emph{Called}, not hooked. ``The hard 20\%.'' \\
\texttt{line_commit}\newline(FUN_00433720) & \texttt{0x33720} &
\ic{__cdecl(void)}; advances the scrollback to the next line. Called after each \texttt{line_append}. \\
\texttt{dispatch_submit}\newline(FUN_0055f800) & \texttt{0x15F800} &
The command dispatcher; fires on Enter, main thread. Reads the typed line from \texttt{DAT_00bf8130}, tokenises, scans the \ic{\{name, handler\}} table, else prints ``Unknown Command''. \emph{Hooked} (capture + route). \\
\texttt{frame_render}\newline(FUN_00575be0) & \texttt{0x175BE0} &
The per-frame terminal renderer. First half = typewriter reveal + \ic{--MORE--} pager; second half = the grid draw loop. \emph{Hooked} as the per-frame pump. \\
\texttt{cell_write}\newline(FUN_00422d30) & \texttt{0x22D30} &
\ic{(ch, color, flag)} $\rightarrow$ writes \ic{grid_chars[cursor]=ch}, \ic{grid_attrs[cursor]=(color<<4)\textbar flag}, advances the cursor; \texttt{\textbackslash n} scrolls. \\
\texttt{key_handler}\newline(FUN_0047eb90) & \texttt{0x7EB90} &
The terminal key handler \ic{(ctx, keycode, down)}. Recovered but \emph{not} used (see Section~\ref{sec:input}). \\
\bottomrule
\end{longtable}

\noindent\textbf{Display globals} (the basis for the emulator):

\begin{longtable}{@{}llp{8.5cm}@{}}
\toprule
\textbf{Global} & \textbf{RVA} & \textbf{Meaning} \\
\midrule
\texttt{DAT_00bf56a0} & \texttt{0x7F56A0} & Char grid. Row-major, stride \texttt{0x50} (80 cols), 27 rows (to offset \texttt{0x86f}). \\
\texttt{DAT_00bf44d8} & \texttt{0x7F44D8} & Attribute grid. Byte \ic{= (color<<4)|flag}; high nibble selects the cell text colour. \\
\texttt{DAT_00bf812c} & \texttt{0x7F812C} & Linear grid cursor \ic{= col + row*0x50}. \\
\texttt{DAT_00bf8120} & \texttt{0x7F8120} & Lines \emph{written} (write head of a 100-row ring). \\
\texttt{DAT_00bf8124} & \texttt{0x7F8124} & Lines \emph{revealed} (reveal head; chases the write head = typewriter). \\
\texttt{DAT_00bf8130} & \texttt{0x7F8130} & The raw typed command line. \\
\bottomrule
\end{longtable}

\section{System architecture}

The injected DLL (\texttt{bo1bridge.dll}, x86, static CRT) is loaded into the game
with a small custom injector (\texttt{tools/inject.cpp}: \texttt{CreateRemoteThread}
+ \texttt{LoadLibrary}). On attach it installs inline hooks with \textbf{MinHook}.
There are two user-facing modes, both toggled by typing a word in the terminal:

\begin{itemize}[nosep]
  \item \texttt{connect} --- \textbf{line-shell mode}: each typed command runs in a
  persistent \texttt{cmd.exe} and the output is rendered with the game's own
  \texttt{line_append}.
  \item \texttt{emu} --- \textbf{emulator mode}: a ConPTY shell drives a VT parser and
  a character grid that is blitted into the engine's display memory; full-screen
  interactive programs (including Claude Code) render here.
\end{itemize}

\subsection{Threading model}

Engine functions must only be touched on the \textbf{main thread}; calling them from
a worker crashes sporadically. The design keeps all engine contact on the per-frame
hook:

\begin{lstlisting}
submit hook ---> [inbound queue] ---> worker/reader thread ---> shell / ConPTY
                                                  |
                                                  v
per-frame hook <--- drain <--- [outbound] / grid model <--- output bytes
   |
   +-- renders ON THE MAIN THREAD ONLY (line_append, or grid blit)
\end{lstlisting}

\noindent
In line-shell mode a worker thread runs the shell and pushes output lines; the frame
hook drips a few lines per frame. In emulator mode a reader thread pumps ConPTY bytes
into the VT parser (under a lock), and the frame hook blits the resulting grid.

\section{Line-shell bridge (\texttt{connect} mode)}

A single persistent \texttt{cmd.exe} runs as a REPL over redirected pipes, so
\texttt{cd} and environment changes persist and there is no per-command process
launch cost. Each command's output is delimited by a unique marker line; reads poll
with a timeout that bounds any main-thread hitch. The submit hook captures the typed
line, runs the original dispatcher with output suppressed (to preserve the engine's
input bookkeeping while hiding the fake built-ins), executes the command on the
worker, and the frame hook streams the result into the scrollback via
\texttt{line_append}. The engine's built-in typewriter reveal and \ic{--MORE--} pager
(renderer ``case 5'') are reused for free; \texttt{Ctrl-C} and \texttt{q} at a pager
break abort a flood of output.

This mode is sufficient for ordinary CLI use (file access, persistent \texttt{cd},
one-shot commands) but cannot host interactive full-screen programs --- which is why
the emulator was built.

\section{Terminal emulator (\texttt{emu} mode)}

This is the core of the project. The pipeline is:

\begin{lstlisting}
ConPTY child (cmd.exe) --> reader thread --> vt::Parser --> grid::Screen (80x27)
        --> emuBlit() paints engine grid each frame --> engine renderer draws it
input: WndProc subclass --> VT bytes --> ConPTY stdin
\end{lstlisting}

\subsection{ConPTY backend (\texttt{src/conpty.\{h,cpp\}})}

A Windows pseudoconsole makes the child believe it is talking to a real terminal: it
emits VT/ANSI and accepts VT input --- exactly the byte stream a terminal emulator
parses. The ConPTY entry points (\texttt{CreatePseudoConsole},
\texttt{ResizePseudoConsole}, \texttt{ClosePseudoConsole}) are resolved dynamically
from \texttt{kernel32} so the build does not depend on SDK header versions. Two pipes
are created (our input $\rightarrow$ child, child output $\rightarrow$ us), the
pseudoconsole is attached to the child through a \texttt{STARTUPINFOEX} process-thread
attribute (\texttt{PROC_THREAD_ATTRIBUTE_PSEUDOCONSOLE}), and the child is launched
with \texttt{CreateProcess}.

\subsection{VT/ANSI parser (\texttt{src/vtparse.\{h,cpp\}})}

A small state machine (ground / escape / CSI) interprets the child's output:

\begin{itemize}[nosep]
  \item \textbf{Cursor + edits:} CUP, ED, EL, DECSTBM scroll region, DECSC/DECRC save
  /restore (\ic{ESC 7/8}, \ic{CSI s/u}), ECH \ic{CSI X}, ICH \ic{CSI @}, DCH \ic{CSI P},
  IL \ic{CSI L}, DL \ic{CSI M}, VPA \ic{CSI d}, RI \ic{ESC M}.
  \item \textbf{Colour (SGR):} the 16 ANSI colours plus 256-colour (\ic{38;5;n}) and
  truecolour (\ic{38;2;r;g;b}), each mapped to nearest ANSI-8 by
  \ic{xterm256ToAnsi8}/\ic{rgbToAnsi8}. Background SGR is consumed but ignored (the
  engine cell stores foreground only).
  \item \textbf{UTF-8:} multi-byte sequences are decoded to a codepoint (one grid cell
  each) and mapped to ASCII: box-drawing $\rightarrow$ \texttt{| - + =}, blocks
  $\rightarrow$ \texttt{\# : .}, braille (spinners) $\rightarrow$ \texttt{*}, smart
  quotes/dashes to ASCII, unknown printables to \texttt{?}. The BO1 font is ASCII-only,
  so this is a readable approximation.
  \item \textbf{Query replies:} DSR (\ic{ESC[6n}) and DA (\ic{ESC[c}) are answered via a
  responder callback that writes back to the pty --- without this, size detection
  stalls (Section~\ref{sec:dsr}).
  \item \textbf{Alternate screen} (\ic{?1049h/?47h}) and a capped scroll-back history.
\end{itemize}

\subsection{Grid model (\texttt{src/grid.\{h,cpp\}})}

\ic{grid::Screen} is an 80$\times$27 array of \ic{\{char, color\}} cells with a cursor,
the current SGR colour, a scroll region, a saved-cursor slot, a scroll-back history
\ic{deque}, a view offset for paging, and an alt-screen flag. Colour is kept in ANSI
space (0--7, bright 8--15); the mapping to the engine palette happens later, at blit
time. A subtle but important member is \ic{wrapPending_}, the deferred-wrap flag
(Section~\ref{sec:decawm}).

\subsection{Input: WndProc subclass (\texttt{src/hooks.cpp})}
\label{sec:input}

Static recon showed the game's keyboard input is \textbf{Windows-message driven} ---
it imports \texttt{RegisterClass}, \texttt{DefWindowProc}, \texttt{CallWindowProc},
the \texttt{Get/Peek/Translate/DispatchMessage} family --- and uses \emph{no}
DirectInput. So rather than hook the engine's awkward internal key handler, the game
window is subclassed (found at runtime via \texttt{EnumWindows} filtered by own PID /
visible / titled). While in emu mode \ic{emuWndProc} takes \texttt{WM_CHAR} characters
and \texttt{WM_KEYDOWN} special keys, translates them to VT input bytes
(arrows $\rightarrow$ \ic{ESC[A..D}, Enter $\rightarrow$ \ic{\textbackslash r}, Backspace
$\rightarrow$ \texttt{0x7f}, PgUp/PgDn, etc.), writes them to the pty, and
\textbf{swallows} the message so the game takes no action (no movement, no ESC pause
menu). \texttt{Ctrl+]} leaves emu mode; PgUp/PgDn page the local scroll-back.

\section{The renderer internals (deep dive)}
\label{sec:renderer}

Understanding \ic{frame_render} (\texttt{FUN_00575be0}) in detail was the key to making
arbitrary content render correctly. It has two halves.

\textbf{First half --- typewriter reveal.} It walks the reveal head
(\texttt{DAT_00bf8124}) toward the write head (\texttt{DAT_00bf8120}), revealing queued
rows over time, and runs the \ic{--MORE--} pager (``case 5''). The emulator
\emph{skips} this half by pinning \ic{reveal == write} each frame, so it never repaints
the line ring over the blitted grid.

\textbf{Second half --- the grid draw loop.} This unconditionally draws all 27 grid
rows. It iterates rows (screen lines 9..35) and, within each row, splits the cells
into \emph{runs} of equal attribute, drawing each run with one call chain:

\begin{lstlisting}
FUN_00575be0 (draw loop)
  -> FUN_005f4a90 (draw one run)
       -> FUN_006d30f0 (font metric), FUN_005cd7b0 (position/justify transform)
       -> FUN_0061c4d0 -> FUN_006d65c0 (the glyph emitter)
  width measured by FUN_005d84f0 -> FUN_006d3110 (parses ^N colour codes)
\end{lstlisting}

Two properties of this path turned out to dictate exactly how the grid must be filled,
and were the root causes of the two biggest bugs.

\subsection{Property A: NUL-terminated rows + a shared vertex budget}
\label{sec:propA}

\ic{FUN_006d65c0} emits a run by reading it as a \textbf{NUL-terminated C-string}
(\ic{strlen} then \ic{memcpy} of that length --- it \emph{ignores} the element count it
is passed) into a \textbf{single per-frame text-vertex buffer}. That buffer has a hard
size guard:

\begin{lstlisting}
size = (strlen(run) + 0x60) & ~3;
if (((buf[4] + buf[2]) - write_offset) - 0x2000 < size) {  // would overflow headroom
    buf[3] = 0;
    return;   // emit NOTHING -- for the rest of the frame
}
\end{lstlisting}

\noindent
So a run with \emph{no} NUL terminator is read far past its row (to the next stray NUL,
roughly 2\,KB away), consuming a huge slice of the shared budget; once the cumulative
budget is exhausted, \emph{every later run that frame silently draws nothing}.

\subsection{Property B: run colour is taken from the boundary cell}
\label{sec:propB}

The draw loop recomputes the draw colour from the \emph{current} cell on every
iteration and emits a run when the attribute changes. At the break, the colour in hand
is the attribute of the cell that \emph{ended} the run --- i.e. the first cell of the
\emph{next} colour segment. The run \ic{[start..boundary]} is therefore drawn entirely
in the \emph{next} segment's colour. Painting each cell its own true colour thus bleeds
every segment onto its neighbour.

How the compensations for both properties were derived appears in
Section~\ref{sec:errcolour} and Section~\ref{sec:errvisible}; the final blit is below.

\subsection{The blit (\texttt{emuBlit})}

\ic{emuBlit} runs on the frame hook and fills the engine grid in two passes:

\begin{enumerate}[nosep]
  \item \textbf{Pass 1 --- glyphs + desired colour.} Copy each grid cell's glyph into
  \texttt{DAT_00bf56a0} and compute the desired engine colour (\ic{ansiToEngine(cell.color)})
  into a parallel buffer. \emph{Spaces inherit the preceding glyph's colour}, so a lone
  one-cell coloured character (e.g. a \texttt{|} border) extends across the following
  space into a two-cell segment the engine can colour (it cannot isolate a one-cell run).
  \item \textbf{Pass 2 --- attributes (filler scheme) + NUL termination.} For each colour
  segment, the interior cells get a ``filler'' attribute equal to the \emph{previous}
  segment's colour (always different from this segment, so the run boundary lands exactly
  on the segment end) and only the segment's \emph{last} cell carries the real colour
  (Property B). Every row is then NUL-terminated after its last glyph so runs stay bounded
  (Property A); a fully blank row becomes all-NUL and is skipped cleanly.
\end{enumerate}

\noindent
After the blit, the frame hook pins \ic{reveal = write} to disable the typewriter half.
The engine palette, calibrated live, is: nibble \texttt{0} = white (the game's native
default), \texttt{2} = black, \texttt{3..8} = red/green/yellow/blue/magenta/cyan,
\texttt{9..F} = white. \ic{ansiToEngine} maps ANSI white to \texttt{0} (no hacks).

\section{Notable and interesting findings}

\begin{enumerate}[label=\textbf{F\arabic*.},leftmargin=*]
  \item \textbf{The terminal display is a real 2D grid in engine memory.} It is not a
  list of scrollback strings --- \ic{cell_write} maintains an 80$\times$27 char grid plus a
  parallel attribute grid at fixed addresses. This is what makes the whole emulator
  possible: we can write our own grid straight into engine memory and let the game's
  existing renderer paint it, with zero graphics-API overlay, preserving the authentic
  CRT look.

  \item \textbf{ASLR is off.} \texttt{BlackOps.exe} always loads at \texttt{0x400000},
  so addresses are stable and globals are reachable as \ic{base + RVA}. This is unusual
  for the era and made both static and dynamic analysis far easier.

  \item \textbf{Run colour is off-by-one by design.} The engine colours a run by the
  cell that \emph{ends} it (Property B, Section~\ref{sec:propB}). This is fine for the
  game's own content (its colour changes mostly fall on spaces) but is hostile to
  arbitrary per-cell colour; it forced the ``filler'' scheme.

  \item \textbf{The renderer shares a capped per-frame vertex buffer and reads rows as
  C-strings.} A run with no NUL terminator is over-read and can exhaust the budget,
  after which the rest of the frame draws nothing (Property A, Section~\ref{sec:propA}).
  This single mechanism explained an entire class of ``the bottom is cut off'' symptoms.

  \item \textbf{Input is plain Win32 messages, no DirectInput.} That made a WndProc
  subclass the clean way to capture keystrokes, far simpler than hooking the engine's
  internal key path.

  \item \textbf{The terminal has a built-in typewriter reveal and a real
  \texttt{-{}-MORE-{}-} pager} (renderer ``case 5''). The line-shell mode reuses both;
  the emulator deliberately disables the reveal by pinning the heads equal.

  \item \textbf{Quake-style \texttt{\textasciicircum{}N} colour codes.} The width/measure
  helper parses \ic{\textasciicircum N} inline colour escapes (0 width) --- a relic of the idTech
  lineage, still live in the terminal's text path.

  \item \textbf{White is the default palette slot (nibble 0), and it is visible.} An
  early misreading held that there was ``no usable white''; in fact nibble 0 is the
  game's native white. The confusion was an artefact of an unrelated bug hiding white
  rows (Section~\ref{sec:errwhite}).

  \item \textbf{A bug-compensation accidentally produced a feature.} The brightness-tracking
  RGB$\rightarrow$ANSI mapping makes Claude Code's pulsing-orange ``thinking'' spinner
  flip between red and yellow as it brightens --- a live shimmer the user specifically
  wanted kept (Section~\ref{sec:shimmer}).

  \item \textbf{32-bit process $\Rightarrow$ WOW64 redirection hides 64-bit tools.}
  Because the game is 32-bit, a naive \texttt{cmd.exe} child sees
  \texttt{System32} redirected to \texttt{SysWOW64}, hiding \texttt{ssh}/\texttt{scp}
  (which exist only in the real \texttt{System32\textbackslash OpenSSH}). The fix is to
  launch the native 64-bit shell via the \texttt{Sysnative} alias
  (Section~\ref{sec:ssh}).
\end{enumerate}

\section{Errors encountered and how they were fixed}

\subsection{First live test froze the game (non-unique signature)}
The \texttt{line_commit} byte signature matched \textbf{ten} locations --- it was the
compiler's generic \ic{(x \% 100) * 0x56} idiom on the line counter, which recurs across
the terminal code. The scanner returned the first match, so the bridge called into the
middle of an unrelated function as \ic{void(void)} and corrupted the stack.
\textbf{Fix:} since ASLR is off, resolve \texttt{line_commit} by \textbf{RVA +
prologue verification} (a guarded read that can never crash on a wrong address) instead
of by pattern; keep unique signatures for the others.

\subsection{Synchronous shell froze the engine thread}
Running the shell synchronously inside the submit hook blocked the engine thread, so
the renderer could not run; output came one command behind, with typewriter and pager
glitches. \textbf{Fix:} move the shell to a \textbf{worker thread} and add a
\textbf{per-frame pump} that drips a few output lines per frame, throttled so the
100-row ring never overruns.

\subsection{PowerShell buffering and padding}
PowerShell block-buffers stdout to a pipe (output arrives a command behind) and pads
lines with trailing spaces. \textbf{Fix:} use \texttt{cmd.exe} as the shell, which
flushes immediately and does not pad.

\subsection{Input-waiting commands hung the session}
A command that reads stdin (\texttt{python}, \texttt{more}, \texttt{sort}, a bare REPL,
or an unrecognised command dropping to a sub-prompt) blocked on the pipe and swallowed
the end-of-output marker, so every later command read the previous command's output.
\textbf{Fix:} append \ic{ <nul} to each command so input-waiters get immediate EOF; on a
read timeout, mark the session desynced and drain the pipe before the next command.

\subsection{cmd echoed piped commands; the marker false-matched}
\texttt{cmd.exe} echoes every piped command line, and \texttt{echo off} hides only the
prompt, not the line --- so the user saw their command reprinted, and the echoed
\ic{echo <marker>} line contained the marker verbatim, causing the reader to latch onto
the echo instead of real output. \textbf{Fix:} send the marker via a caret-escaped
command so cmd echoes it \emph{with} the caret (no false match) while its output drops
the caret (real match); then strip the echoed command line and any line containing the
marker stem.

\subsection{ConPTY never signalled EOF on child exit}
The pseudoconsole keeps the output pipe's write end open, so \ic{read()} blocked forever
after the child exited. \textbf{Fix:} do blocking reads on a dedicated reader thread and
call \texttt{ClosePseudoConsole} on teardown to unblock it.

\subsection{ConPTY child leaked to the parent console}
A console-subsystem test harness let the child attach to the \emph{parent} console and
write there instead of into the pty (the pipe got only the init burst).
\textbf{Fix:} make the harness a \textbf{GUI-subsystem} process (as the game is), so the
child has no parent console to attach to and must use the pty.

\subsection{Visible height miscalibrated (20 vs 27 rows)}
An early estimate put the visible area at ${\sim}20$ rows; a temporary row-number overlay
showed rows 0--26 are all visible. At 20/24 rows, Claude Code's bottom-anchored input box
fell below the rendered area. \textbf{Fix:} set \ic{kEmuRows = 27} (the full grid).

\subsection{Terminal size detection stalled ${\sim}15$\,s}
\label{sec:dsr}
Claude Code rendered inconsistently with a ${\sim}15$\,s delay before first output. The
tell was the delay: TUIs detect terminal size by emitting a query (\ic{ESC[6n} cursor
report, often after \ic{ESC[999;999H}, and \ic{ESC[c} device attributes) and
\emph{blocking} for the reply. The parser never replied, so the app timed out and
fell back to a mis-detected size. \textbf{Fix:} a responder callback writes replies back
to the pty --- DSR $\rightarrow$ \ic{ESC[<row>;<col>R} (which, after the app's
\ic{ESC[999;999H}, clamps to the real 80$\times$27), DA $\rightarrow$ \ic{ESC[?1;2c}.

\subsection{Eager line wrap caused drift and duplication (DECAWM)}
\label{sec:decawm}
The grid wrapped eagerly when the 80th column was written, producing an extra scroll;
over a session this shifted the screen up, duplicated content, and scrolled the input
box out of view (the ``can't see my input / doesn't scroll far enough'' symptom).
\textbf{Fix:} implement \textbf{deferred wrap} (DECAWM) --- writing the last column sets
\ic{wrapPending_}; the wrap happens only when the \emph{next} glyph arrives, and every
cursor move / CR / LF / scroll-region / clear resets the flag. Verified offline by
replaying captured pty logs through a host harness (\texttt{tools/feed_log.cpp}).

\subsection{The big one: the bottom rows were culled}
\label{sec:errvisible}
The most stubborn bug: Claude Code's lower rows --- especially the input box ---
did not render, even though the engine grid memory provably held the correct content
and the draw loop iterates all 27 rows. A key diagnostic (a row-number overlay) showed
only rows that previously held dense content rendered; sparse rows did not. A uniform
``colour test'' that filled \emph{every} row with dense text rendered all 27 rows.

The cause was \textbf{Property A} (Section~\ref{sec:propA}): the blit filled rows with
spaces and \emph{never wrote a NUL}, so each colour run was read far past its row to the
next stray NUL. Claude's mixed-colour rows produce many runs each; the top rows drained
the shared vertex budget and every later run --- the bottom rows --- was dropped. The
uniform test rendered because uniform rows are one run each and fit the budget. Earlier
theories (the write head sitting too low; forcing \ic{written = shown = 27}) had
\emph{no} effect and were removed. \textbf{Fix:} NUL-terminate each grid row after its
last glyph, bounding every run to its own row.

\subsection{Colour bled between segments}
\label{sec:errcolour}
With visibility fixed, colours were scrambled: borders and text took their neighbour's
colour, with only boundary cells correct. This was \textbf{Property B}
(Section~\ref{sec:propB}) --- runs coloured by their ending cell. \textbf{Fix:} the
\textbf{filler scheme}: real colour only on a segment's last cell, interior cells set to
the previous segment's colour so run boundaries align with segment ends; plus space
colour propagation so one-cell borders become colourable two-cell segments. A lone
one-cell colour segment with no adjacent space is the only residual imperfection ---
rare in practice.

\subsection{``No white'' misdiagnosis}
\label{sec:errwhite}
While chasing the visibility bug, body text (which is white) was invisible, leading to
the wrong conclusion that the engine had no visible white, and a temporary
``white $\rightarrow$ green'' hack. Once the real cause (Property A) was fixed, nibble 0
was confirmed to be visible white --- the game's native default. \textbf{Fix:} drop the
hack; map ANSI white $\rightarrow$ engine nibble 0 directly.

\subsection{Logo and box were different colours}
Claude's bright cream logo mapped to ``white'' while the box's deeper orange mapped to
red, because the original \ic{rgbToAnsi8} used a flat per-channel \ic{>127} threshold that
split one hue across buckets. \textbf{Fix:} split on \emph{saturation} --- low-saturation
colours (the cream logo) map to white, while saturated colours pick a hue.

\subsection{The spinner stopped shimmering}
\label{sec:shimmer}
The saturation fix above, in its first form, biased saturated oranges toward solid red,
which flattened the ``thinking'' spinner's red$\leftrightarrow$yellow pulse. That pulse
came from brightness: as the spinner brightens, its green channel crosses the threshold,
flipping the quantised colour. \textbf{Fix:} keep the saturation gate (so the cream logo
stays white) but use the per-channel \ic{>127} threshold for \emph{saturated} colours, so
the brightening orange flips red$\leftrightarrow$yellow again while the static box orange
stays red.

\subsection{\texttt{ssh} ``not recognized''}
\label{sec:ssh}
\texttt{ssh} failed in the emulator shell. The first hypothesis --- a stale inherited
PATH --- was wrong: rebuilding PATH from the registry actually \emph{broke} \texttt{claude}
(it dropped the user PATH) and did not fix \texttt{ssh}. The real cause is
\textbf{WOW64 file-system redirection}: the game is 32-bit, so launching \texttt{cmd.exe}
yields the 32-bit cmd, inside which \texttt{C:\textbackslash Windows\textbackslash System32}
is redirected to \texttt{SysWOW64}; \texttt{ssh}/\texttt{scp} live only in the real
64-bit \texttt{System32\textbackslash OpenSSH}, so they are invisible. \textbf{Fix:}
launch the native 64-bit cmd via the \texttt{Sysnative} alias
(\ic{C:\textbackslash Windows\textbackslash Sysnative\textbackslash cmd.exe}), which a 32-bit process can use to reach the real
\texttt{System32}; leave PATH untouched. The same change starts the shell in
\ic{\%USERPROFILE\%} instead of the game's install folder (a new \ic{cwd} argument on the
ConPTY session).

\subsection{Development friction}
Two recurring, non-product issues: (a) the DLL could not be rebuilt while the game held
it loaded (linker \texttt{LNK1104}) --- the game must be closed before each build; and
(b) chaining a \texttt{git commit} with a here-string under PowerShell failed, so commits
were made through a POSIX shell.

\section{Result and current state}

Claude Code runs and renders correctly inside the Black Ops terminal: the welcome box,
colours, box-drawing borders (as ASCII), the live red$\leftrightarrow$yellow ``thinking''
spinner, scroll-back paging, and keyboard input all work; \texttt{ssh}, \texttt{python},
and ordinary CLI tools run from a native 64-bit shell started in the user's home
directory. All debug scaffolding has been removed. The reverse-engineering record lives
in \texttt{notes/offsets.md} (the curated function/global table) and
\texttt{notes/re-log.md} (the chronological log), backed by the raw decompiles in
\texttt{notes/} and the re-runnable Ghidra scripts in \texttt{ghidra_scripts/}.

The project demonstrates, end to end, that a modern interactive TUI can be hosted inside
a 2010 game's hardcoded easter-egg terminal --- with no source, no symbols, and no
graphics-API overlay --- purely by reverse engineering the engine's own text renderer
and feeding it from a pseudoconsole.

\end{document}
